\section{Formalization}
\label{formalization}

We formalize our basic construction in Figure \ref{fig:OTS-PC}. We make use of an auxiliary function to create Taproot addresses out of scripts:

\begin{enumerate}
    \item \texttt{TaprootAddress}(list-of-scripts): a P2TR address with a tree containing the given scripts with unspendable NUMS internal key.
    
\end{enumerate}

We make use of an auxiliary functions create verify onetime signatures against a list of public keys: 
\begin{enumerate}  
    \item \texttt{\opc{OT-CSIGV}}(list of OT-public keys): a script that reads from the stack an encoded input message $M$ and one or more OT signatures and verifies the signatures of $M$ against a list of OT public keys. Aborts the script if any of the verifications fail. Leaves the value $M$ on top of the stack. 
\end{enumerate}


We define a transaction as \tuple{list-of-inputs, list-of-outputs}. 
Each input is defined as a prevout structure (\tuple{TxID,index}). However, since we have named each output, we use the name instead. An output is defined as \tuple{amount, address}.
We omit transaction versions and nSequence, and lockTime values for simplicity. One exception must be kept in mind by implementers: \opc{CHECKSEQUENCEVERIFY} is enforced only if the spending transaction has version $\geq 2$ and its input sets nSequence to a relative lock at least as large as the one in the script. A transaction that omits either silently loses the timelock, and with it the dispute window. We additionally use:

\begin{itemize}
\item "-" to indicate that we don't care about this field.

\item \opc{Covenant-Check($i$)} to indicate the signing by all parties involved (Alice and Bob in our case). This can be accomplished by checking 2 signatures using a \opc{CHECKSIGVERIFY} followed by a \opc{CHECKSIG}, or by checking one signature aggregated for both parties using MuSig2 with a single \opc{CHECKSIG}. If present, an index ($i$) indicates which taproot path is being signed. It is written last in every script below, for the reason given at the end of this list.

\item The MuSig2 option requires Taproot. A MuSig2 aggregate is a Schnorr key, and \opc{CHECKSIG} verifies Schnorr signatures only in Taproot; segwit v0 verifies ECDSA over a compressed key and rejects a 32-byte one outright. So an output whose script contains a MuSig2 \opc{Covenant-Check} cannot be P2WSH, and every covenant output below is Taproot for that reason. Taking the P2WSH option and the MuSig2 option together yields an output that can never be spent, and the two are described independently, so nothing in either description looks wrong.

\item \opc{CovAddress} to indicate an address that is an aggregated for Alice and Bob using MuSig2.

\item \opc{CSIG}$(x)$ is defined as the script that pushes the public key $x$, followed by \opc{CHECKSIG}. This is the \emph{terminal} form of a signature check.

\item \opc{CSIGV}$(x)$ is defined as the script that pushes the public key $x$, followed by \opc{CHECKSIGVERIFY}. This is the non-terminal form, for a signature check that is not the last in its script.

\item \opc{CSEQV}$(x)$ is defined as the script that pushes value $x$,  and executes  \opc{CHECKSEQUENCEVERIFY} followed by \opc{DROP}.

\item \opc{CHASHV}$(x)$ is defined as the script that executes \texttt{HASH160}, then pushes the hash value $x$, and executes \texttt{EQUALVERIFY}. Note the order: the opcode hashes the preimage taken from the stack and compares the result against $x$.

\item \opc{CVALV}$(x)$ is defined as the script that pushes the value $x$, and executes \texttt{EQUALVERIFY}.

\item A witness script must finish with exactly one true value on the stack. Each of \opc{CSEQV}, \opc{CHASHV}, \opc{CVALV} and the \opc{LESSTHAN} \opc{VERIFY} of the punishment consumes its operands and pushes nothing, so a script assembled only from them terminates with an \emph{empty} stack, and the output it guards can never be spent. Nothing about that fails noisily: the address is well formed, the spending transaction is well formed, and the defect appears only when a spend is attempted and rejected. Every script below therefore places its signature check last and in the terminal \opc{CSIG} form, so that the value \opc{CHECKSIG} pushes is what satisfies the rule. The ordering is not cosmetic; a reader transcribing these scripts with the checks in any other order will produce unspendable outputs.

\end{itemize}

To facilitate reading, we also embed scripts in outputs as pre-P2SH, while in practice all outputs would be Taproot and specify script hashes, while scripts are revealed in the witness stack in the spending transaction.

We assign names to outputs, so we can refer to the outputs by name instead of by transaction id/index.

We omit the definitions of WTA and WTB disputes since they are similar to Alice and Bob disputes respectively, only paying a small fee to a  watchtower address out of the total funds paid.

The balance $I_{BAL}$ represents the total balance deposited in the payment channel. Balances $A_{BAL}(i)$ and $B_{BAL}(i)$ represents the balances of Alice and Bob after state update $i$, not including the amounts that may be locked in HTLCs nor the amounts committed to be paid in transaction fees.

We define $\epsilon$ as the amount in bitcoins that is the minimum necessary to make the largest transaction in our graph standard. The standard transaction rules dictate that a certain fee-rate must be paid by transaction inputs, even when child transactions are paying for parent fees using CPFP. The \txname{Setup} transaction moves $k\epsilon$ bitcoins to each exit output, where $k$ is the worst-case depth of transactions below it. In the basic OTS-PC construction of this appendix the worst case is $k=3$ for an exit output (\txname{CommitExit}, \txname{AssertExitState}, \txname{FinalizeExit}) and $k=2$ for a dispute output (\txname{CommitPunishX}, \txname{PunishX}). The privacy variants of section \ref{watchtowers} insert a transaction into one of these chains and raise the corresponding $k$ by one; the budget each requires is given there.

The scheme uses 3 addresses defined during setup: AliceAddress, BobAddress and SharedAddress. AliceAddress and BobAddress are the addresses where Alice and Bob will receive the funds after unilateral exit. The SharedAddress is used only to let Alice or Bob bump transaction fees using the Child-Pay-For-Parent (CPFP) method. The only transactions that need this output are \txname{AliceAssertExitState} and \txname{BobAssertExitState}. The remainder transactions can be accelerated by Alice or Bob using the package relay method. This requires creating a new transaction \txname{Forward} to transfer the newly received (but unconfirmed) bitcoins to a new owned address, paying higher-than-normal transaction fees. With package relay nodes/miners accept and relay packages of related transactions instead of each transaction in isolation. In our case, it makes miners include all (unconfirmed) parent transactions of \txname{Forward}.

Each owner has its own exit path, rooted at its own \ioname{Alice\_Exit\_out} or \ioname{Bob\_Exit\_out} output. Both paths are defined symmetrically, so below we give Alice's path in full and only note Bob's where it differs. Although every transaction is signed by both owners, the signature on a path's first transaction is delivered in one direction only: Bob gives Alice his signature on \txname{AliceCommitExit($i$)} and never receives hers, and vice versa. Each owner can therefore publish its own path and neither can publish the other's.

\txname{Setup} = \tuple{
    \tuple{ \ioname{Funds\_From\_Alice\_in}, \ioname{Funds\_From\_Bob\_in}},
    \tuple{ \ioname{Initial\_Funds\_out}, \ioname{Alice\_Exit\_out}, \ioname{Bob\_Exit\_out},
    \ioname{Bob\_Disputes\_out}, \ioname{Alice\_Disputes\_out}, 
    \ioname{WTB\_Disputes\_out}, \ioname{WTA\_Disputes\_out} }
}

\ioname{Funds\_From\_Alice\_in} =\tuple{-, -}

\ioname{Funds\_From\_Bob\_in}   =\tuple{-, -} 

\ioname{Initial\_Funds\_out} = \tuple{ $I_{BAL}$ ,  \opc{CovAddress} }

\ioname{Alice\_Exit\_out} = \tuple{ $3\epsilon$, \opc{CovAddress} } 

\ioname{Bob\_Exit\_out} = \tuple{ $3\epsilon$, \opc{CovAddress} } 

\ioname{Bob\_Disputes\_out} = \tuple{ $2\epsilon$, \opc{CovAddress} } 

\ioname{Alice\_Disputes\_out} = \tuple{ $2\epsilon$, \opc{CovAddress} } 

\ioname{WTB\_Disputes\_out} = \tuple{ $2\epsilon$,  \opc{CovAddress} } 

\ioname{WTA\_Disputes\_out} = \tuple{ $2\epsilon$,  \opc{CovAddress} } 
 
\txname{AliceCommitExit($i$)} = \tuple{
    \tuple{ \ioname{Alice\_Exit\_out}} ,
    \tuple{ \ioname{AliceCommitExit\_out}}
}

\ioname{AliceCommitExit\_out} = \tuple{ $2\epsilon$, 
 \texttt{TaprootAddress}( \tuple{ 
      \opc{OT-CSIGV}( $K^A_{pub}$ ) \opc{CVALV($i$)}
      \opc{Covenant-Check}
      })
 }

\txname{BobCommitExit($i$)} and \ioname{BobCommitExit\_out} are identical with \ioname{Bob\_Exit\_out} and $K^B_{pub}$ in place of \ioname{Alice\_Exit\_out} and $K^A_{pub}$.

\txname{AliceAssertExitState($i$)} = \tuple{
    \tuple{ \ioname{AliceCommitExit\_out}} ,
    \tuple{ \ioname{AliceReady\_out},\ioname{AliceCPFP\_out}}
}

\ioname{AliceReady\_out} = \tuple{ $\epsilon$, {
 \texttt{TaprootAddress}(\tuple{ 
 \tuple{ 
 \textbf{[1]} 
  \opc{CSEQV}( $T$ ) 
  \opc{Covenant-Check(1)} },
 \tuple{ 
 \textbf{[2]} 
   \opc{CHASHV($H_a$)} 
   \opc{CHASHV($H_b$)} 
   \opc{Covenant-Check(2)} 
  }
 })
 }
}

\ioname{AliceCPFP\_out}= \tuple{ $\epsilon$, SharedAddress }

\txname{AliceFinalizeExit($i$)} = \tuple{
    % inputs
    \tuple{ \ioname{Initial\_Funds\_out}, \ioname{AliceReady\_out}},
    %outputs
    \tuple{
        \tuple{ $A_{BAL}(i)$, AliceAddress }
        \tuple{ $B_{BAL}(i)$, BobAddress }
    }
}

\txname{BobAssertExitState($i$)}, \ioname{BobReady\_out}, \ioname{BobCPFP\_out} and \txname{BobFinalizeExit($i$)} mirror the above. Note that \txname{AliceFinalizeExit($i$)} and \txname{BobFinalizeExit($i$)} pay the same amounts to the same addresses, and both spend \ioname{Initial\_Funds\_out}: only one of them can ever confirm, and it does not matter which.

%%%%%%%%%%%%%%%%%%%%%%%%%%%%%
\txname{CommitPunishAlice($i$)} = \tuple{
    % inputs
    \tuple{ \ioname{Bob\_Disputes\_out}},
    %outputs
    \tuple{ \ioname{PunishAlice\_out} }
}
    
\ioname{PunishAlice\_out} = \tuple{ $\epsilon$, \texttt{TaprootAddress}({
          \tuple{ 
           \opc{OT-CSIGV}( $K^A_{pub}$ )  
           \tuple{i} 
           \opc{LESSTHAN} 
           \opc{VERIFY} 
           \opc{Covenant-Check} } 
        }
}

\txname{PunishAlice($i$)} = \tuple{
    % inputs
    \tuple{ \ioname{Initial\_Funds\_out}, \ioname{PunishAlice\_out}},
    %outputs
    \tuple{
        \tuple{ $I_{BAL}$, BobAddress }
    }
}

%%%%%%%%%%%%%%%%%%%%%%%%%%%
\txname{CommitPunishBob($i$)} = \tuple{
    % inputs
    \tuple{ \ioname{Alice\_Disputes\_out}},
    %outputs
    \tuple{ \ioname{PunishBob\_out} }
}
    
\ioname{PunishBob\_out} = \tuple{ $\epsilon$, \texttt{TaprootAddress}({
          \tuple{ 
           \opc{OT-CSIGV}( $K^B_{pub}$ )  
           \tuple{i} 
           \opc{LESSTHAN} 
           \opc{VERIFY} 
           \opc{Covenant-Check} } 
        }
}

\txname{PunishBob($i$)} = \tuple{
    % inputs
    \tuple{ \ioname{Initial\_Funds\_out}, \ioname{PunishBob\_out}},
    %outputs
    \tuple{
        \tuple{ $I_{BAL}$, AliceAddress }
    }
}

%%%%%%%%%%%%%%%%%%%%%%%%%%%%%%%%%
